% Gesamttest «Exponential- und Logarithmusfunktionen» — Quelle des PDFs (python3 scripts/build-lp-pdf.py).
% Musterlösung und Raster: bewertungspaket.tex. Alle Werte mit python3 nachgerechnet (04.10.2026).
% Kein \lt/\gt — pdfLaTeX kennt sie nicht, < und > tun es.
\documentclass[11pt]{article}
\usepackage{../lp-druck}
\renewcommand{\lpname}{Exponential- und Logarithmusfunktionen}
\renewcommand{\lpteilgebiet}{3.4}
\renewcommand{\lpbereich}{Schwerpunktfach}
\renewcommand{\lpfuss}{Gesamttest · 25 Punkte}
\begin{document}
\lpkopf{Gesamttest}{%
  \vspace{4pt}\noindent\begin{tabularx}{\linewidth}{@{}X X@{}}
  Code (kein Name): \hrulefill & Punkte: \hrulefill\ / 25
  \end{tabularx}}

\begin{lpinfo}{Anleitung}
\textbf{Zeit:} rund 30 Minuten. \textbf{Hilfsmittel:} kein Taschenrechner — alle vier Kompetenzen des Teilgebiets 3.4 tragen im Lehrplan den Vermerk «auch ohne Hilfsmittel». Lineal und Bleistift sind erlaubt.
Schreib zu jeder Aufgabe den \textbf{Rechenweg} auf — bewertet wird der Weg.
Danach: Lösung scannen oder fotografieren (ohne Namen und Standort) und mit dem \emph{Bewertungspaket} bewerten lassen.
\end{lpinfo}

\lpteil{Teil A · Exponentialfunktion, Wachstum und Zerfall}{Kapitel 1 und 2 · 13 Punkte}

\begin{lpaufgabe}{G1}{Zwei Kurven skizzieren}{3 P}
\begin{minipage}[t]{0.52\linewidth}
Skizziere $y = 3^{x}$ und $y = \left(\tfrac13\right)^{x}$ für $-2 \le x \le 2$ in das Koordinatensystem. Gib den gemeinsamen Punkt, die Asymptote und die Wertemenge an.
\lpplatz{6}
\end{minipage}\hfill
\begin{minipage}[t]{0.42\linewidth}\vspace{0pt}\centering
\begin{tikzpicture}
\begin{axis}[width=6.4cm,height=6.4cm,axis lines=middle,xmin=-2.4,xmax=2.4,ymin=-1,ymax=10,
  xtick={-2,-1,1,2},ytick={2,4,6,8},minor ytick={1,3,5,7,9},minor xtick={-1.5,-0.5,0.5,1.5},grid=both,grid style={lplinie},tick label style={font=\scriptsize},
  xlabel=$x$,ylabel=$y$,
  yticklabel style={fill=white,inner sep=1pt,xshift=-2pt},xticklabel style={fill=white,inner sep=1pt,yshift=-2pt}]
\end{axis}
\end{tikzpicture}
\end{minipage}
\end{lpaufgabe}

\begin{lpaufgabe}{G2}{Die Basis bestimmen}{3 P}
Bestimme die Basis $a$ der Exponentialfunktion $y = a^{x}$, deren Graph durch den Punkt geht:
\begin{enumerate}[label=(\alph*),leftmargin=*,itemsep=1pt]
\item $P(-3 \mid 8)$
\item $Q\left(\tfrac32 \,\middle|\, 27\right)$
\end{enumerate}
\lpplatz{6}
\end{lpaufgabe}

\begin{lpaufgabe}{G3}{Eine Population wächst}{4 P}
Eine Tierpopulation wird jedes Jahr gezählt ($t$ in Jahren):
\begin{center}\begin{tabular}{l|ccc}
$t$ & $0$ & $1$ & $2$\\\hline
$N(t)$ & $500$ & $600$ & $720$
\end{tabular}\end{center}
\begin{enumerate}[label=(\alph*),leftmargin=*,itemsep=1pt]
\item Zeig, dass das Wachstum exponentiell und nicht linear ist, und stell das Modell $N(t)$ auf.
\item Berechne $N(3)$.
\item Um wie viel Prozent wächst die Population pro Jahr?
\end{enumerate}
\lpplatz{7}
\end{lpaufgabe}

\begin{lpaufgabe}{G4}{Ein Medikament wird abgebaut}{3 P}
Von einem Medikament sind $240$ mg im Blut; die Halbwertszeit beträgt $6$ Stunden.
\begin{enumerate}[label=(\alph*),leftmargin=*,itemsep=1pt]
\item Stell das Modell $m(t)$ auf ($t$ in Stunden) und berechne $m(18)$.
\item Nach wie vielen Stunden sind noch $15$ mg übrig?
\end{enumerate}
\lpplatz{6}
\end{lpaufgabe}

\lpteil{Teil B · e-Funktion, Sättigung, Logarithmusfunktion}{Kapitel 3 bis 5 · 12 Punkte}

\begin{lpaufgabe}{G5}{Basiswechsel}{3 P}
\begin{enumerate}[label=(\alph*),leftmargin=*,itemsep=1pt]
\item Schreib $25^{x}$ in der Form $5^{b x}$.
\item Schreib $\left(\tfrac19\right)^{x}$ in der Form $3^{b x}$.
\item Schreib $7^{x}$ in der Form $e^{b x}$.
\end{enumerate}
\lpplatz{5}
\end{lpaufgabe}

\begin{lpaufgabe}{G6}{Wasser wird erwärmt}{3 P}
Die Temperatur (in °C) von Wasser in einem Wasserkocher folgt dem Modell $f(t) = 90 - 70\,e^{-k t}$ ($t$ in Minuten, $k > 0$).
\begin{enumerate}[label=(\alph*),leftmargin=*,itemsep=1pt]
\item Gib Starttemperatur und Sättigungswert an.
\item Steigt oder fällt die Kurve? Wie heisst ihre waagrechte Asymptote?
\item Der Abstand zum Sättigungswert halbiert sich alle $5$ Minuten. Welche Temperatur hat das Wasser nach $10$ Minuten?
\end{enumerate}
\lpplatz{6}
\end{lpaufgabe}

\begin{lpaufgabe}{G7}{Die Umkehrfunktion}{4 P}
Gegeben ist $f(x) = 2^{x} - 3$.
\begin{enumerate}[label=(\alph*),leftmargin=*,itemsep=1pt]
\item Bestimme die Umkehrfunktion $f^{-1}$ mit Rechenweg.
\item Gib Definitionsmenge und Asymptote von $f^{-1}$ an.
\item Berechne die Nullstelle von $f^{-1}$.
\end{enumerate}
\lpplatz{7}
\end{lpaufgabe}

\begin{lpaufgabe}{G8}{Die Logarithmuskurve durch Spiegeln}{2 P}
\begin{minipage}[t]{0.52\linewidth}
Abgebildet sind $y = 2^{x}$ und die Gerade $y = x$ (gestrichelt). Zeichne den Graphen von $y = \log_2 x$ ein, indem du die Kurve an $y = x$ spiegelst. Markiere die Nullstelle und den Punkt mit $x = 4$ und gib beide als Koordinaten an.
\lpplatz{4}
\end{minipage}\hfill
\begin{minipage}[t]{0.42\linewidth}\vspace{0pt}\centering
\begin{tikzpicture}
\begin{axis}[width=6.4cm,height=6.4cm,axis lines=middle,xmin=-2.4,xmax=5.4,ymin=-2.4,ymax=5.4,
  xtick={-2,-1,1,2,3,4,5},ytick={-2,-1,1,2,3,4,5},grid=both,grid style={lplinie},tick label style={font=\scriptsize},
  xlabel=$x$,ylabel=$y$,
  yticklabel style={fill=white,inner sep=1pt,xshift=-2pt},xticklabel style={fill=white,inner sep=1pt,yshift=-2pt}]
\addplot[lpblau,thick,domain=-2.4:2.4,samples=80]{2^x};
\addplot[black!55,dashed,domain=-2.4:5.4]{x};
\end{axis}
\end{tikzpicture}
\end{minipage}
\end{lpaufgabe}
\end{document}
